\section{Conclusion}
We introduced FANTOM, a unified framework for multi-regime MTS that jointly infers \textit{(i)} the number of regimes, \textit{(ii)} their boundaries, and \textit{(iii)} their corresponding causal DAG, under either non-Gaussian or heteroscedastic noises. Under mild assumptions in causal discovery, we prove identifiability of the temporal heteroscedastic causal models in the stationary case and show that the number of regimes, their indices, and their graphs are identifiable (up to permutation) in the non-stationary setting. Extensive experiments on synthetic and real-world data show consistent gains over strong baselines. FANTOM offers a principled means to uncover regime-specific causal dynamics, enhancing regime detection, and causal discovery with potential applications in various domains such as finance, climate science, and neuroscience.
